% Generated draft: requires mathematical and source review.
\documentclass[11pt]{article}
\usepackage[margin=25mm]{geometry}
\usepackage{fontspec}
\setmainfont{Latin Modern Roman}
\usepackage{amsmath,amssymb,mathtools,unicode-math}
\setmathfont{Latin Modern Math}
\usepackage{enumitem,needspace,xurl,hyperref}
\setlength{\parindent}{0pt}
\setlength{\parskip}{5pt}
\setlength{\emergencystretch}{3em}
\newcommand{\N}{\mathbb{N}}
\newcommand{\Z}{\mathbb{Z}}
\newcommand{\Q}{\mathbb{Q}}
\newcommand{\R}{\mathbb{R}}
\newcommand{\C}{\mathbb{C}}
\newcommand{\F}{\mathbb{F}}
\newcommand{\A}{\mathbb{A}}
\newcommand{\PP}{\mathbb{P}}
\newcommand{\E}{\mathbb{E}}
\newcommand{\eps}{\varepsilon}
\newcommand{\dd}{\,\mathrm{d}}
\begin{document}
\section*{UnsolvedMath ID: 800187}
\textbf{Source problem code:} catalogue entry
\subsection{Definitions and mathematical statement}
For a positive integer \(k\) and a positive integer \(m\), let the generalized divisor function be
\[
d_k(m)=\#\{(a_1,\ldots,a_k)\in\Z_{\ge1}^{\,k}:a_1\cdots a_k=m\}.
\]
Let \(r\ge3\) be an integer, let \(k_1,\ldots,k_r\) be positive integers, and let \(h_1,\ldots,h_r\) be pairwise distinct positive integers. For real \(x\ge1\), define
\[
S_{\mathbf{k},\mathbf{h}}(x)=\sum_{\substack{n\in\Z\,;\ 1\le n\le x}}\prod_{j=1}^{r}d_{k_j}(n+h_j).
\]
The problem is to formulate an explicit conjectural asymptotic formula for \(S_{\mathbf{k},\mathbf{h}}(x)\) as \(x\to\infty\). In particular, seek an explicit leading expression \(A_{\mathbf{k},\mathbf{h}}(x)\) such that one conjectures
\[
\frac{S_{\mathbf{k},\mathbf{h}}(x)}{A_{\mathbf{k},\mathbf{h}}(x)}\longrightarrow1.
\]
Here the parameters are held fixed when taking the limit; the supplied source does not specify uniformity in them, the desired error term, or the extent of any lower-order expansion. The source writes \((k_1,\ldots,k_r)\in\N^r\); the formulation above uses the positive-integer convention. If zero is intended to belong to \(\N\), the source does not specify the meaning of \(d_0\).

The source motivates this problem by analogy with conjectural Hardy--Littlewood asymptotics for
\[
\sum_{\substack{n\in\Z\,;\ 1\le n\le x}}\Lambda(n+h_1)\cdots\Lambda(n+h_r),
\]
where \(\Lambda(m)=\log p\) if \(m=p^a\) for a prime \(p\) and an integer \(a\ge1\), and \(\Lambda(m)=0\) otherwise. This analogy does not supply a formula for the divisor correlation. The source also states that Ng--Thom and Tao used different probabilistic arguments to derive conjectures for \(r=2\), but supplies no specific references for those works.
\subsection{Short English statement}
For every fixed choice of at least three pairwise distinct positive shifts and positive divisor-function orders, conjecture an explicit asymptotic formula for the sum of the corresponding shifted divisor-function products. The source compares this task with Hardy--Littlewood prime correlations and mentions probabilistic conjectures for two shifts. This is a statement-only draft requiring human mathematical review.
\subsection{Sources}
\begin{itemize}
\item Local UnsolvedMath registry, permanent ID \(800187\), source code \texttt{catalogue entry}, in the AMR Open Problem Lists collection. The supplied record attributes the problem to Hamieh, Kadiri, and Ng, with proposed year 2024, and labels it \texttt{open}. This label is supplied registry metadata, not an independently verified mathematical or literature assessment. The record was created and updated on 2026-10-16; upstream revision: \texttt{d371ee917ac9647fbb06cd6fc529deec6f3dda82}.
\item Hamieh, Kadiri, and Ng, \emph{Moments of L-functions Problem List} (2024), Problem 39, under Moments of the Riemann zeta function. Supplied source URL: \url{https://arxiv.org/abs/2405.17800}. The background records access on 2026-10-16 and extraction by \texttt{arxiv\_latex}; these are supplied provenance statements, not claims of access or verification for this draft.
\item The supplied statement invokes the Hardy--Littlewood conjectures as an analogy and attributes probabilistic conjectures for \(r=2\) to Ng--Thom and Tao. It provides no bibliographic details or links for these contextual references.
\end{itemize}
\end{document}
